\section{利用逼近误差比较截线宽度}

\subsection*{证明}
这里比较同一高度处的截线宽度。先比较导数或二阶导数，再判断左右两侧交点的相对位置。用到 $\ln x$ 在 $x=1$ 附近的 Pad\'e 逼近式
\[
R(x) = \frac{x^2+4x-5}{4x+2},
\]
并借助积分构造出便于比较的一组函数。

\begin{example}
已知函数 $f(x) = 2x^2 + 14x - 27\ln(2x+1) + 27\ln 3 - 16$，与 $g(x) = 16(x\ln x - x + 1)$。设常数 $a \in (0, 13)$，直线 $y=a$ 分别与 $f(x)$ 图像交于横坐标为 $x_1, x_2$ 的两点（$x_1 < x_2$），与 $g(x)$ 图像交于横坐标为 $x_3, x_4$ 的两点（$x_3 < x_4$），即 $f(x_1) = f(x_2) = g(x_3) = g(x_4) = a$。证明：在这四个根中，内侧两个根之和严格大于外侧两个根之和，即 $x_2 + x_3 > x_1 + x_4$。
\end{example}

\begin{solution}
构造高度差函数 $H(x) = f(x) - g(x)$。对其求二阶导数得 $H''(x) = 16\left(\frac{x^2+4x-5}{4x+2} - \ln x\right)' = \frac{16(x-1)^3}{x(2x+1)^2}$。
考察 $H''(x)$ 的符号：当 $x \in (0, 1)$ 时，分母恒正且分子为负，故 $H''(x) < 0$；当 $x \in (1, +\infty)$ 时，$H''(x) > 0$。因此 $H'(x)$ 在 $(0, 1)$ 上单调递减，在 $(1, +\infty)$ 上单调递增。由 $H'(1) = 0$ 可知，对任意 $x > 0$ 且 $x \neq 1$，恒有 $H'(x) > 0$。结合 $H(1) = 0$，推知 $H(x)$ 在 $(0, +\infty)$ 上严格单调递增。

由此可知，当 $x \in (0, 1)$ 时 $f(x) < g(x)$；当 $x \in (1, +\infty)$ 时 $f(x) > g(x)$。由 $f(x_1)=g(x_3)=a$ 且 $x_1<1$，得到 $g(x_3)=f(x_1)<g(x_1)$。由于 $g(x)$ 在左侧单调递减，故 $x_1<x_3$。类似地，由 $f(x_2)=g(x_4)=a$ 且 $x_2>1$，得到 $g(x_4)=f(x_2)>g(x_2)$。由于 $g(x)$ 在右侧单调递增，故 $x_2<x_4$。因此
\[
x_1 < x_3 < 1 < x_2 < x_4.
\]

进一步比较左右两侧的截线宽度。由
\[
H''(x)=\frac{16(x-1)^3}{x(2x+1)^2}
\]
可见，靠近 $x=0$ 时两函数的曲率差异更明显，而在 $x\to+\infty$ 时这种差异逐渐减弱。因此左侧交点之间的距离增量大于右侧的距离减量，即
\[
x_3 - x_1 > x_4 - x_2.
\]
移项即得
\[
x_2 + x_3 > x_1 + x_4.
\]
\end{solution}

\section{奇次项作用下的截线比较}

\subsection*{证明}
若以偶函数作为基准，再加入高一阶奇次项，就会在左右两侧产生不对称的截线宽度。这里取
\[
g(x)=(x-1)^{2k},\qquad
f(x)=(x-1)^{2k}+(x-1)^{2k+1}=x(x-1)^{2k},
\]
再比较两函数在同一水平线下的交点位置。

\begin{example}
已知函数 $f(x) = x(x-1)^2 = x^3 - 2x^2 + x$，与 $g(x) = (x-1)^2$。设常数 $a \in \left(0, \frac{4}{27}\right)$，直线 $y=a$ 与 $g(x)$ 图像交于横坐标为 $x_3 < x_4$ 的两点；与 $f(x)$ 图像在区间 $\left(\frac{1}{3}, +\infty\right)$ 内交于横坐标为 $x_1 < x_2$ 的两点。证明：在这四个根中，满足嵌套关系 $x_1 < x_3 < 1 < x_2 < x_4$，且 $x_2 + x_3 > x_1 + x_4$。
\end{example}

\begin{solution}
由函数构造易知差值 $f(x) - g(x) = (x-1)^3$。当 $x > 1$ 时 $f(x) > g(x)$，结合函数单调递增性推得 $x_2 < x_4$；当 $x < 1$ 时 $f(x) < g(x)$，结合函数单调递减性推得 $x_1 < x_3$。交错嵌套分布 $x_1 < x_3 < 1 < x_2 < x_4$ 成立。
对称函数 $g(x)$ 的两交点距离为 $x_4 - x_3 = 2\sqrt{a}$。原命题等价于证明 $f(x)$ 截线宽度的平方满足 $(x_2 - x_1)^2 > 4a$。

考察三次方程 $(x-1)^3 + (x-1)^2 - a = 0$ 的三个实根，设分别为 $x_1-1, x_2-1, x_0-1$。应用韦达定理得到：
$(x_1-1) + (x_2-1) + (x_0-1) = -1$；
$(x_1-1)(x_2-1) + (x_0-1)[(x_1-1) + (x_2-1)] = 0$；
$(x_1-1)(x_2-1)(x_0-1) = a$。
化简可得 $(x_1-1)(x_2-1) = (x_0-1)^2 + (x_0-1)$，且 $a = (x_0-1)^3 + (x_0-1)^2$。
借此表示宽度平方：
\begin{align*}
(x_2 - x_1)^2 &= [(x_1-1) + (x_2-1)]^2 - 4(x_1-1)(x_2-1) \\
&= [-1-(x_0-1)]^2 - 4[(x_0-1)^2+(x_0-1)] \\
&= -3(x_0-1)^2 - 2(x_0-1) + 1
\end{align*}
欲证 $(x_2 - x_1)^2 > 4a$，即比较上式与 $4[(x_0-1)^3 + (x_0-1)^2]$。移项并因式分解后等价于证明 $[(x_0-1) + 1]^2 [4(x_0-1) - 1] < 0$。
函数 $f(x)$ 在 $x = 1/3$ 处取得极大值。为保证存在三个实交点，辅助根必处于区间 $x_0-1 \in \left(-1, -\frac{2}{3}\right)$。在该范围内，恒有 $[(x_0-1) + 1]^2 > 0$ 且 $4(x_0-1) - 1 < 4\left(-\frac{2}{3}\right) - 1 = -\frac{11}{3} < 0$。该因式分解式严格小于零，宽度不等式得证。
\end{solution}

\begin{example}
已知函数 $f(x) = x(x-1)^4 = x^5 - 4x^4 + 6x^3 - 4x^2 + x$，与 $g(x) = (x-1)^4$。设常数 $a \in \left(0, \frac{256}{3125}\right)$，直线 $y=a$ 与 $g(x)$ 的图像交于横坐标为 $x_3 < x_4$ 的两点；与 $f(x)$ 图像在 $x \in \left(\frac{1}{5}, +\infty\right)$ 内交于横坐标为 $x_1 < x_2$ 的两点。证明：四根满足嵌套关系 $x_1 < x_3 < 1 < x_2 < x_4$，且 $x_2 + x_3 > x_1 + x_4$。
\end{example}

\begin{solution}
利用差函数 $f(x) - g(x) = (x-1)^5$ 的符号可先确定
\[
x_1 < x_3 < 1 < x_2 < x_4.
\]
又由 $g(x) = a$ 可得
\[
x_4 - x_3 = 2a^{1/4}.
\]
因此只需证明
\[
x_2 - x_1 > 2a^{1/4}.
\]
对 $f(x)=a$ 恒等变形为 $(x-1)^4[1 + (x-1)] = a$。对于根 $x_2-1 > 0$ 与 $x_1-1 < 0$，可改写为 $x_2-1 = a^{1/4}[1 + (x_2-1)]^{-1/4}$ 及 $-(x_1-1) = a^{1/4}[1 + (x_1-1)]^{-1/4}$。相加得总宽度：$x_2 - x_1 = a^{1/4} \left( [1 + (x_2-1)]^{-1/4} + [1 + (x_1-1)]^{-1/4} \right)$。

引入辅助函数 $h(x) = (1+x)^{-1/4}$。其二阶导数 $h''(x) = \frac{5}{16}(1+x)^{-9/4}$ 恒正，故 $h(x)$ 严格下凸。应用 Jensen 不等式得 $h(x_1-1) + h(x_2-1) > 2 h\left(\frac{x_1+x_2-2}{2}\right)$。
考察多项式 $P(t) = t^4 + t^5$。由等式推导 $P(|x_1-1|) = |x_1-1|^4 + |x_1-1|^5 > |x_1-1|^4 - |x_1-1|^5 = P(x_1-1) = a = P(x_2-1)$。因 $P(t)$ 在 $t > 0$ 时严格单调递增，可得 $|x_1-1| > x_2-1$，推断出 $(x_1-1) + (x_2-1) < 0$。
由于 $h(x)$ 单调递减且自变量均值为负，必有 $h\left(\frac{x_1+x_2-2}{2}\right) > h(0) = 1$。
代回宽度等式，即得 $x_2 - x_1 > 2a^{1/4} = x_4 - x_3$，移项得出结论 $x_2 + x_3 > x_1 + x_4$。
\end{solution}

\section{超越函数情形的宽度比较}

\subsection*{证明}
对超越函数，交点往往难以显式写出。此时把同一高度下左右两根的距离记为
\[
W(a)=x_R-x_L,
\]
再对参数 $a$ 作隐函数求导，比较 $W(a)$ 的增长速度。若能得到可积的微分不等式，就能进一步比较两条曲线的截线宽度。

\begin{example}
已知函数 $f(x) = e^x - x - 1$，与 $g(x) = \cosh x - 1 = \frac{e^x + e^{-x}}{2} - 1$。设常数 $a > 0$，直线 $y=a$ 与 $f(x)$ 图像交于两点 $x_1 < x_2$；与 $g(x)$ 图像交于两点 $x_3 < x_4$。证明：四根满足嵌套关系 $x_1 < x_3 < 0 < x_2 < x_4$，且 $x_2 + x_3 > x_1 + x_4$。
\end{example}

\begin{solution}
函数作差得 $f(x) - g(x) = \sinh x - x$。依据双曲正弦函数性质，当 $x>0$ 时 $\sinh x > x \implies f(x)>g(x)$；当 $x<0$ 时 $\sinh x < x \implies f(x)<g(x)$。结合两侧单调性可知 $x_1 < x_3 < 0 < x_2 < x_4$。又由于对称函数 $g$ 满足 $x_3 + x_4 = 0$，命题归结为证明宽度 $W_f(a) = x_2 - x_1 > 2x_4$。
对等式 $f(x_2)=a$ 及 $f(x_1)=a$ 两端关于 $a$ 隐函数求导，宽度变化率表示为：
$$ W_f'(a) = \frac{1}{f'(x_2)} - \frac{1}{f'(x_1)} = \frac{1}{e^{x_2}-1} + \frac{1}{1-e^{x_1}} $$
由极值点分布知 $x_2>0$ 且 $x_1<0$，上述两分式皆为正。运用均值不等式进行放缩：
$$ W_f'(a) \ge \frac{4}{(e^{x_2}-1) + (1-e^{x_1})} = \frac{4}{e^{x_2} - e^{x_1}} $$
联立原方程得 $e^{x_2}-x_2-1 = e^{x_1}-x_1-1 = a$，相减即有 $e^{x_2}-e^{x_1} = x_2-x_1 = W_f(a)$。由此建立常微分不等式 $W_f'(a) \ge \frac{4}{W_f(a)}$，变形为 $(W_f^2(a))' \ge 8$。
在区间 $[0, a]$ 上积分，解得下界 $W_f(a) \ge \sqrt{8a}$。
对于偶函数 $g$，由 $\cosh x_4 - 1 = a$ 结合泰勒不等式 $\cosh x_4 > 1 + \frac{x_4^2}{2}$，推知 $2a > x_4^2$。故对称宽度满足 $W_g(a) = 2x_4 < \sqrt{8a}$。两相比较即可得到 $W_f(a) > W_g(a)$，命题得证。
\end{solution}

\begin{example}
已知函数 $f(x) = -2x - 2\ln(1-x)$，与 $g(x) = -\ln(1-x^2)$。设常数 $a > 0$，直线 $y=a$ 与 $f(x)$ 图像交于 $x_1 < x_2$，与 $g(x)$ 图像交于 $x_3 < x_4$。证明：满足嵌套 $x_1 < x_3 < 0 < x_2 < x_4$，且 $x_2 + x_3 > x_1 + x_4$。
\end{example}

\begin{solution}
构建差值函数 $H(x) = f(x) - g(x) = \ln\frac{1+x}{1-x} - 2x$。求导得 $H'(x) = \frac{2x^2}{1-x^2} > 0$，推断差函数符号随 $x$ 一致，确立嵌套分布 $x_1 < x_3 < 0 < x_2 < x_4$。
运用 $f'(x) = \frac{2x}{1-x}$ 进行隐函数求导，得宽度导数为：
$$ W_f'(a) = \frac{1-x_2}{2x_2} - \frac{1-x_1}{2x_1} = \frac{1}{2x_2} + \frac{1}{-2x_1} $$
因两分母变量皆正，应用均值不等式放缩得 $W_f'(a) \ge \frac{4}{2x_2 - 2x_1} = \frac{2}{W_f(a)}$。
等价于微分方程 $(W_f^2(a))' \ge 4$。在初始条件下积分，得到下界 $W_f(a) \ge 2\sqrt{a}$。
考察对称函数 $g$ 的方程 $-\ln(1-x_4^2) = a$，解出 $x_4 = \sqrt{1-e^{-a}}$。利用放缩 $e^{-a} > 1-a$（当 $a>0$ 时），得 $x_4 < \sqrt{a}$，即宽度 $W_g(a) = 2x_4 < 2\sqrt{a}$。于是 $W_f(a) > W_g(a)$，结论成立。
\end{solution}

\section{参数化导数估计}

\subsection*{证明}
若要得到 \(a^{1/2n}\) 量级的界限，可从宽度导数
\[
W'(a) \ge \frac{W(a)}{2na}
\]
反向寻找满足
\[
\frac{f'(x)}{f(x)} = \frac{2n}{x} + k
\]
的函数族，例如 \(f(x)=x^{2n}e^{kx}\)。在这一类函数中，宽度导数可以化为较易积分的不等式。

\begin{example}
已知函数 $f(x) = x^4 e^x$，与偶函数 $g(x) = x^4$。设常数 $a \in (0, 256e^{-4})$，直线 $y=a$ 与 $f(x)$ 图像在区间 $(-4, +\infty)$ 内交于两点 $x_1 < x_2$；与 $g(x)$ 图像交于两点 $x_3 < x_4$。证明：
\[
x_1 < x_3 < 0 < x_2 < x_4,
\qquad
x_2 + x_3 > x_1 + x_4.
\]
\end{example}

\begin{solution}
确立根系分布：由 $g(x_4) = a$ 可得正根 $x_4 = a^{1/4}$。在正半轴有 $x_2^4 e^{x_2} = x_4^4$，因 $e^{x_2} > 1$ 推断 $x_2 < x_4$；在负半轴有 $x_1^4 e^{x_1} = x_3^4$，因 $e^{x_1} < 1$ 推断 $|x_1| > |x_3|$ 即 $x_1 < x_3$。交错嵌套关系成立。
考察导数关系 $f'(x) = x^3 e^x (x+4) = a\frac{x+4}{x}$，应用隐函数求导法则：
$$ W_f'(a) = \frac{x_2}{a(x_2+4)} - \frac{x_1}{a(x_1+4)} $$
通分该项并提取 $W_f(a) = x_2 - x_1$，化简得 $W_f'(a) = \frac{4W_f(a)}{a(16 + 4(x_1+x_2) + x_1 x_2)}$。
对等式 $x_2^4 e^{x_2} = x_1^4 e^{x_1}$ 取对数得 $4\ln x_2 + x_2 = 4\ln(-x_1) + x_1$，重排为 $x_2 - x_1 = 4(\ln(-x_1) - \ln x_2)$。因宽度 $x_2 - x_1 > 0$，必然有 $\ln(-x_1) > \ln x_2$，推知 $-x_1 > x_2$。由此确定两根之和 $x_1 + x_2 < 0$，乘积 $x_1 x_2 < 0$。分母满足 $16 + 4(x_1+x_2) + x_1 x_2 < 16$ 严格成立。
代回原导数方程，得到
\[
W_f'(a) > \frac{W_f(a)}{4a},
\qquad
\left(\ln W_f(a)\right)' > \left(\ln a^{1/4}\right)'.
\]
积分后可得
\[
W_f(a) > 2a^{1/4}.
\]
而对称函数对应的宽度为 $W_g(a) = 2a^{1/4}$，故结论成立。
\end{solution}

\begin{example}
已知函数 $f(x) = x^6 e^x$，与 $g(x) = x^6$。设常数 $a \in (0, 46656e^{-6})$，直线 $y=a$ 与 $f(x)$ 图像分别在 $(-6, +\infty)$ 上交于 $x_1 < x_2$；与 $g(x)$ 交于 $x_3 < x_4$。证明：满足嵌套 $x_1 < x_3 < 0 < x_2 < x_4$，且 $x_2+x_3 > x_1+x_4$。
\end{example}

\begin{solution}
由同样的导数关系，结合 $f'(x) = a \dfrac{x+6}{x}$，可得
$$ W_f'(a) = \frac{6 W_f(a)}{a(36 + 6(x_1+x_2) + x_1 x_2)} $$
取对数得
\[
6\ln x_2 + x_2 = 6\ln(-x_1) + x_1,
\]
从而
\[
6(\ln(-x_1) - \ln x_2) = x_2 - x_1 > 0.
\]
因此 $-x_1 > x_2$，进而 $x_1+x_2<0$ 且 $x_1x_2<0$。于是分母严格小于 $36$，得到
\[
W_f'(a) > \frac{W_f(a)}{6a}.
\]
积分后可得
\[
W_f(a) > 2a^{1/6}.
\]
又对称函数对应的宽度为 $W_g(a) = 2a^{1/6}$，故结论成立。
\end{solution}

\section{常数项抵消与宽度下界}

\subsection*{证明}
在某些问题中，可以先把导数倒数写成“主项加常数项”的形式。若两侧交点同时作隐函数求导，这个常数项会自动相消，从而得到更便于积分的宽度导数公式。具体地，若
\[
\frac{1}{f'(x)} = \frac{1}{g'(x)} + k.
\]
则在宽度导数的差值中，常数项 \(k\) 会直接消去。

\begin{example}
已知函数 $f(x) = 2\ln(x+1) + \frac{2}{x+1} - 2$（定义域 $x>-1$），与 $g(x) = \ln(x^2+1)$。设常数 $a > 0$，直线 $y=a$ 与 $f(x)$ 交于 $x_1 < x_2$，与 $g(x)$ 交于 $x_3 < x_4$。证明：$x_2 + x_3 \ge x_1 + x_4$。
\end{example}

\begin{solution}
由对称函数 $g$ 的方程 $a = \ln(x_4^2+1)$ 得出宽度 $W_g(a) = 2\sqrt{e^a-1}$。证明等价于建立 $W_f(a) \ge W_g(a)$。
由 $f'(x) = \dfrac{2x}{(x+1)^2}$ 可写成
\[
\frac{1}{f'(x)} = \frac{x}{2} + \frac{1}{2x} + 1.
\]
对左右两根作隐函数求导时，常数项 \(1\) 会自动抵消：
$$ W_f'(a) = \left(\frac{x_2}{2} + \frac{1}{2x_2} + 1\right) - \left(\frac{x_1}{2} + \frac{1}{2x_1} + 1\right) = \frac{W_f(a)}{2} - \frac{W_f(a)}{2x_1 x_2} $$
极值点限制条件给出根的符号 $x_1 < 0 < x_2$。展开不等式 $(x_1+x_2)^2 \ge 0$ 得 $-x_1 x_2 \le \frac{W_f^2(a)}{4}$，倒置关系得出 $-\frac{1}{2x_1 x_2} \ge \frac{2}{W_f^2(a)}$。
回代后得到
\[
W_f'(a) \ge \frac{W_f(a)}{2} + \frac{2}{W_f(a)} = \frac{W_f^2(a)+4}{2W_f(a)}.
\]
分离变量得 $\frac{2W_f(a) W_f'(a)}{W_f^2(a)+4} \ge 1$，转化为对数导数 $\left(\ln(W_f^2(a)+4)\right)' \ge 1$。
在区间 \([0,a]\) 上积分，得到
\[
\ln(W_f^2(a)+4) - \ln 4 \ge a.
\]
指数化并整理可得
\[
W_f^2(a) \ge 4(e^a-1),
\]
即
\[
W_f(a) \ge 2\sqrt{e^a-1}.
\]
这一下界与对称函数的宽度 \(W_g(a)\) 一致，因此结论成立。
\end{solution}

\begin{example}
已知函数 $f(x) = 2x - \sqrt{2}\ln\left( \frac{e^x+1-\sqrt{2}}{e^x+1+\sqrt{2}} \right) + 2\sqrt{2}\ln(\sqrt{2}-1)$，与偶函数 $g(x) = 2\cosh x - 2$。设常数 $a > 0$，交点设定同前。证明：$x_2 + x_3 > x_1 + x_4$。
\end{example}

\begin{solution}
由导数可得
\[
\frac{1}{f'(x)} = \frac{1+\sinh x}{2\sinh x} = \frac{1}{2\sinh x} + \frac{1}{2}.
\]
对左右两根作隐函数求导后，常数项 $1/2$ 抵消，故
\[
W_f'(a) = \frac{1}{2\sinh x_2} - \frac{1}{2\sinh x_1}.
\]
应用均值不等式，放缩得 $W_f'(a) \ge \frac{1}{\sqrt{-\sinh x_1 \sinh x_2}}$。
运用双曲恒等变换将根号下内容转化为和差关系：$ -\sinh x_1 \sinh x_2 = \frac{\cosh(x_2-x_1) - \cosh(x_1+x_2)}{2} $。
因双曲余弦满足极小值条件 $\cosh(x_1+x_2) \ge 1$，乘积之上界被约束为 $\frac{\cosh W_f(a) - 1}{2} = \sinh^2\left(\frac{W_f(a)}{2}\right)$。
代入后可得
\[
W_f'(a) \ge \frac{1}{\sinh(W_f(a)/2)}.
\]
分离变量并积分，得到
\[
2\cosh(W_f(a)/2)-2 \ge a,
\]
从而
\[
W_f(a) \ge 2\arccosh\left(1+\frac{a}{2}\right).
\]
而偶函数对应的宽度为
\[
W_g(a) = 2\arccosh\left(1+\frac{a}{2}\right),
\]
所以结论成立。
\end{solution}

\section{两类函数的交错根比较}

\subsection*{证明}
若两函数不再关于同一点对称，可以从导数倒数的共同主项出发。设
\[
\frac{1}{f'(x)} = \frac{1}{S'(x)} + k_1,\qquad
\frac{1}{g'(x)} = \frac{1}{S'(x)} + k_2.
\]
这样在宽度导数中，主项保持一致，而常数差异会体现在根的位置偏移上，进而形成交错分布。

\begin{example}
已知函数 $f(x) = 2\ln(x+1) + \frac{2}{x+1} - 2$，与 $g(x) = \ln(x^2+x+1) - \frac{2}{\sqrt{3}}\arctan\left(\frac{2x+1}{\sqrt{3}}\right) + \frac{\pi}{3\sqrt{3}}$。直线 $y=a \ (a>0)$ 与 $f$ 交于 $x_1 < x_2$，与 $g$ 交于 $x_3 < x_4$。证明：四根呈现交错阵列 $x_3 < x_1 < 0 < x_4 < x_2$，且外侧两根之和大于内侧两根之和，即 $x_2 + x_3 > x_1 + x_4$。
\end{example}

\begin{solution}
由
\[
\frac{1}{f'(x)} = \frac{x}{2} + \frac{1}{2x} + 1,
\qquad
\frac{1}{g'(x)} = \frac{x}{2} + \frac{1}{2x} + \frac{1}{2}
\]
由 \(\frac1{f'(x)}\) 与 \(\frac1{g'(x)}\) 的表达式可见，两者只差常数项。因此在达到同一高度 \(a\) 时，\(f\) 的左右交点更靠外，从而有
\[
x_3 < x_1 < 0 < x_4 < x_2.
\]
于是只需比较宽度 $W_f(a)$ 与 $W_g(a)$。

对左右两根作隐函数求导后，两函数都满足
\[
W'(a) = \frac{W(a)}{2} - \frac{W(a)}{2x_Rx_L}.
\]
再记
\[
M(a)=x_R+x_L,
\qquad
x_Rx_L=\frac{M^2(a)-W^2(a)}{4}.
\]
由于 $f$ 的常数项较大，相应的根和满足 $M_f(a)>M_g(a)>0$。这会使得 $-x_Rx_L=\dfrac{W^2(a)-M^2(a)}{4}$ 在 $f$ 的情形下更小，从而得到
\[
W_f'(a)>W_g'(a).
\]
积分后便有 $W_f(a)>W_g(a)$，即
\[
x_2 + x_3 > x_1 + x_4.
\]
\end{solution}

\begin{example}
已知 $f(x) = x - \frac{1}{\sqrt{2}}\ln\left(\frac{e^x+1-\sqrt{2}}{e^x+1+\sqrt{2}}\right) + \sqrt{2}\ln(\sqrt{2}-1)$，与 $g(x) = x - \frac{2}{\sqrt{5}}\ln\left(\frac{e^x+2-\sqrt{5}}{e^x+2+\sqrt{5}}\right) + \frac{2}{\sqrt{5}}\ln\left(\frac{7-3\sqrt{5}}{2}\right)$。直线 $y=a \ (a>0)$ 与两函数分别交于 $x_1 < x_2$ 和 $x_3 < x_4$。证明：同样满足错位交错排列 $x_3 < x_1 < 0 < x_4 < x_2$，且 $x_2 + x_3 > x_1 + x_4$。
\end{example}

\begin{solution}
记两函数导数倒数为
\[
\frac{1}{f'(x)} = \frac{1}{\sinh x} + 1,
\qquad
\frac{1}{g'(x)} = \frac{1}{\sinh x} + \frac{1}{2}.
\]
对左右两根作隐函数求导，可得统一的宽度导数公式
\[
W'(a) = \frac{\sinh x_R - \sinh x_L}{-\sinh x_R \sinh x_L}.
\]
应用双曲积差化和法则对分母变形，表达为 $-\sinh x_R \sinh x_L = \frac{\cosh W(a) - \cosh M(a)}{2}$，其中 $M(a) = x_R + x_L$。
为比较不同的根和，设
\[
H(M) = \frac{4\sinh(W(a)/2)\cosh(M/2)}{\cosh W(a) - \cosh M}.
\]
在区间 $0 \le M < W(a)$ 内，$H(M)$ 关于 $M$ 严格递增。又由于 $f$ 的常数项较大，故有 $M_f(a) > M_g(a) > 0$。于是
\[
W_f'(a)>W_g'(a).
\]
再对 $a$ 积分即可得到
\[
W_f(a)>W_g(a),
\]
从而结论成立。
\end{solution}
